\documentclass[11pt]{article}
\usepackage{amsmath,amssymb,amsthm,amsfonts}
\usepackage{geometry}
\usepackage{hyperref}
\usepackage{enumitem}
\usepackage{listings}
\usepackage{color}

\geometry{margin=1in}
\hypersetup{
  colorlinks=true,
  linkcolor=blue,
  citecolor=blue,
  urlcolor=blue
}

\definecolor{codegray}{rgb}{0.4,0.4,0.4}
\definecolor{codepurple}{rgb}{0.5,0,0.5}
\lstdefinestyle{mystyle}{
  basicstyle=\ttfamily\small,
  keywordstyle=\color{blue},
  commentstyle=\color{codegray},
  stringstyle=\color{codepurple},
  breaklines=true,
  showstringspaces=false
}
\lstset{style=mystyle}

\title{From Codon-Contrast to IFMD-Governed PIRTM:\\
A Formal Architecture for Spectral Genomics and ZK Certification}
\author{Citizen Gardens \\ \\ The Foundation of Multiplicity}
\date{April 3, 2026}

\begin{document}
\maketitle

\begin{abstract}
This report compiles the developments from a multi-stage design thread:
starting from a 6-bit codon encoding and length-64 Walsh--Hadamard Transform (WHT),
through a series of corrections that led to a fully governed spectral
architecture grounded in the Infinity-Factorization Multiplicity Dynamics (IFMD)
note and integrated with PIRTM's ACE runtime and certified ZK commitments.
We formalize the mathematical objects, map them to concrete circuit and
governance artifacts, and specify an implementation-ready 7-day plan that
meets both mathematical rigor and patent-law constraints under 35 U.S.C. \S112.
\end{abstract}
\newpage
\tableofcontents
\newpage
\section{Executive Summary}

\subsection{Context and Goal}
The starting point is a codon-level representation:
a 6-bit biophysical mapping
\[
  [GC_1, PUR_1, GC_2, PUR_2, GC_3, PUR_3]
\]
combined with an in-place length-64 Walsh--Hadamard transform for genomic
data encoding (``Codon\_Contrast'' specification).
The goal is to transform this into:

\begin{itemize}[noitemsep]
  \item a mathematically sound, cross-domain spectral framework,
  \item a zero-knowledge (ZK) commitment substrate over codon populations, and
  \item a PIRTM/IFMD-governed contraction certificate (``C6'') for the CRMF stack.
\end{itemize}

\subsection{Key Outcomes}

\begin{enumerate}[label=(\alph*),noitemsep]
  \item \textbf{Correct Spectral Object:} We replaced the initial
  character-function accumulation with a two-step pipeline:
  histogram accumulation followed by FWHT, producing a single
  canonical spectral object $\hat{h} \in \mathbb{Z}^{64}$.

  \item \textbf{Epistatic Interpretation:} We showed that
  $\hat{h}(\zeta)/N$ is exactly the $|S_\zeta|$-th order
  population-level epistatic coefficient over the 6-channel fiber,
  consistent with recent WHT epistasis literature.

  \item \textbf{ZK Substrate:} We defined a Poseidon2-based commitment
  to $\hat{h}$ and a per-coefficient predicate
  $\hat{h}(\zeta_{\text{trait}})/N \ge \theta_{\text{trait}}$,
  yielding biologically interpretable ZK proofs.

  \item \textbf{Contraction Certificate (C6):} We replaced CLT-based
  tautologies with IFMD-governed contraction via a data-dependent
  intensity $X_t$, an SNR threshold $\tau_{\min}$, and a WAC
  (Windowed Average Contraction) holdout protocol.

  \item \textbf{Governed Retry Mechanism:} We introduced a
  \emph{retry\_nonce} field in the base parameter bundle, converting
  dataset ``retries'' into an explicit ACE budget burn metric.

  \item \textbf{Circuit Realization:} We arrived at a concrete Groth16
  circuit with $\approx 5087$ constraints, using a
  Poseidon2 sponge with $(t=9,r=8)$ and a minimal trait predicate
  encoding (1 multiplication + 1 range check).
\end{enumerate}

\section{Mathematical Architecture}

\subsection{Biophysical Fiber Space and Codon Encoding}

Let the codon-position base be
\[
  \mathcal{B}_3 = \{1,2,3\},
\]
and per-position fiber
\[
  V_i = \text{span}_{\mathbb{Z}_2}\{ e_{GC_i}, e_{PUR_i} \} \cong \mathbb{Z}_2^{2}, \quad i=1,2,3.
\]
The total feature space is
\[
  \mathcal{F} = V_1 \oplus V_2 \oplus V_3 \cong \mathbb{Z}_2^{6}.
\]
Each codon is encoded as $\omega \in \{0,1\}^{6}$ corresponding to
$(GC_1,PUR_1,GC_2,PUR_2,GC_3,PUR_3)$.

\subsection{Prime-Channel Assignment and Spectral Weights}

Assign distinct primes to channels $k \in \{1,\dots,6\}$:
\[
  p_k \in \{2,3,5,7,11,13\}.
\]
For a spectral index $\xi \in \{0,1\}^6$ with support
$S_\xi = \{k : \xi_k = 1\}$, define the primorial weight:
\[
  w(\xi) = \prod_{k \in S_\xi} p_k, \quad w(\mathbf{0}) = 1.
\]
Define the prime-weighted spectral metric:
\[
  d_{\mathbb{P}}(\hat{\sigma}, \hat{\tau}) =
  \left( \sum_{\xi \in \{0,1\}^6}
    w(\xi)\,|\hat{\sigma}(\xi) - \hat{\tau}(\xi)|^2
  \right)^{1/2}.
\]

\subsection{Histogram and Spectral Projection}

Given a population of codon encodings
$\{\omega_1,\dots,\omega_N\} \subset \{0,1\}^6$,
define the histogram
\[
  h[\xi] = \sum_{n=1}^N \mathbf{1}[\omega_n = \xi], \quad \xi \in \{0,1\}^6,
\]
with $\sum_\xi h[\xi] = N$.

The Walsh--Hadamard transform of $h$ is
\[
  \hat{h}[\zeta] =
  \sum_{\xi \in \{0,1\}^6} h[\xi] (-1)^{\langle \zeta, \xi \rangle},
  \quad \zeta \in \{0,1\}^6,
\]
computed in-place in $O(64 \log 64) = O(384)$ additions.

\subsection{Epistatic Decomposition}

Let $\hat{\mu}(\xi) = h[\xi]/N$ be the empirical codon usage distribution,
and $\hat{\mu}_{\text{WHT}}(\zeta) = \hat{h}(\zeta)/N$ its WHT.
Then
\[
  \hat{\mu}_{\text{WHT}}(\zeta)
  = \sum_{\xi \in \{0,1\}^6} \hat{\mu}(\xi) (-1)^{\langle \zeta, \xi \rangle}.
\]
For $|\zeta| = r$, $\hat{\mu}_{\text{WHT}}(\zeta)$ is the $r$-th order
epistatic coefficient over channels $S_\zeta$, background averaged over
all other channels. The vector $\hat{h}/N$ simultaneously encodes
all orders $0$ through $6$.

\subsection{Spectral Entropy and Sparsity}

Define the normalized spectral mass
\[
  \hat{p}^{(n)}(\zeta) =
  \frac{|\hat{h}^{(n)}(\zeta)|}{\|\hat{h}^{(n)}\|_1},
\]
and spectral entropy
\[
  H_s^{(n)} = -\sum_\zeta \hat{p}^{(n)}(\zeta)
  \log \hat{p}^{(n)}(\zeta).
\]
Convergence is declared at $N_0$ when
\[
  |H_s^{(N_0+1)} - H_s^{(N_0)}| < \delta,
\]
and the $\varepsilon$-support is locked as
\[
  \text{supp}_\varepsilon(\hat{h}^{(N_0)}) =
  \left\{
    \zeta : |\hat{h}^{(N_0)}(\zeta)| >
    \varepsilon \|\hat{h}^{(N_0)}\|_1
  \right\},
\]
with $K = |\text{supp}_\varepsilon| \ll 64$.

\section{ZK Commitment and Predicate}

\subsection{Commitment to Spectral Codon Fingerprint}

The canonical ZK commitment targets the integer vector
$\hat{h} \in \mathbb{Z}^{64}$:
\[
  \mathcal{C}_{\hat{h}} = \text{Poseidon2}(\hat{h})
  \in \mathbb{GF}(p)^\lambda.
\]
A CAS registration commits
\[
  \mathcal{C}_\Gamma = \text{Poseidon2}(\Gamma_d \,\Vert\, \Theta_{C6}),
\]
where:
\begin{itemize}[noitemsep]
  \item $\Gamma_d$ is the channel assignment specification (domain $d$),
  \item $\Theta_{C6}$ is the spectral governance parameter bundle.
\end{itemize}

\subsection{Per-Coefficient Trait Predicate}

For a trait index $\zeta_{\text{trait}} \in \{0,1\}^6$ and threshold
$\theta_{\text{trait}}$, the predicate is
\[
  \frac{\hat{h}(\zeta_{\text{trait}})}{N} \ge \theta_{\text{trait}}.
\]
This is encoded in R1CS as:
\[
  \hat{h}(\zeta_{\text{trait}}) \ge
  \theta_{\text{trait}} \cdot \hat{h}[\mathbf{0}],
\]
using the fact that $\hat{h}[\mathbf{0}] = N$.

\subsection{ZK Relation}

The ZK witness relation is:
\[
  \mathcal{R}_{ZK} =
  \left\{
    (\mathcal{C}_{\hat{h}}, \mathcal{C}_\Gamma, N_{\min}, \zeta_{\text{trait}}, \theta_{\text{trait}})
    \;\middle|\;
    \begin{aligned}
      &\exists\, (\hat{h}, N):\\
      &\text{Poseidon2}(\hat{h}) = \mathcal{C}_{\hat{h}},\\
      &\text{Poseidon2}(\Gamma_d \Vert \Theta_{C6}) = \mathcal{C}_\Gamma,\\
      &\hat{h}(\zeta_{\text{trait}})/N \ge \theta_{\text{trait}},\\
      &N \ge N_{\min}.
    \end{aligned}
  \right\}.
\]

\section{IFMD-Governed Contraction (C6)}

\subsection{Parameter Bundle and Retry Mechanism}

The IFMD-aligned C6 parameter bundle is:
\[
  \Theta_{C6} = (\varepsilon, \delta, N_0, K, M, \beta, \tau_{\min},
  \alpha_M, \text{shuffle\_seed}, \text{domain}, \text{retry\_nonce}).
\]
Here:
\begin{itemize}[noitemsep]
  \item $\varepsilon$ is the ACE gap parameter,
  \item $\delta$ is the spectral entropy tolerance,
  \item $N_0$ is the convergence onset index,
  \item $K$ is the spectral sparsity bound,
  \item $M$ is the WAC window length,
  \item $\beta$ is the holdout split fraction,
  \item $\tau_{\min}$ is the SNR guard (Moonshine Governor),
  \item $\alpha_M$ is the Dirichlet decay exponent,
  \item $\text{shuffle\_seed}$ fixes the gene permutation deterministically,
  \item $\text{retry\_nonce}$ counts governed dataset retries.
\end{itemize}

\subsection{Contraction Intensity from IFMD}

IFMD defines the ACE budget radius:
\[
  R_t = 1 - \varepsilon_t - X_t,
\]
where $X_t \in [0,1)$ is a telemetry-derived contraction intensity.
In the spectral setting, we set
\[
  X_n = \frac{\sum_{\zeta \notin \text{supp}_\varepsilon}
    |\hat{h}^{(n)}(\zeta)|^2}
  {\sum_{\zeta} |\hat{h}^{(n)}(\zeta)|^2},
\]
the off-support spectral variance ratio at accumulation step $n$.
The C6 gate requires:
\[
  \text{support\_size} = K \ll 64, \quad
  \text{SNR} \ge \tau_{\min}, \quad
  X_n \le 1 - \varepsilon
\]
over the WAC window.

\subsection{WAC and Holdout Split}

A windowed average contraction condition (IFMD §3.2) is:
\[
  \|K_t \cdots K_{t+M-1}\| \le \rho < 1, \quad \forall t,
\]
for some $\rho < 1$. In practice:
\begin{itemize}[noitemsep]
  \item Use a holdout split: first $\beta N$ samples estimate
  the attractor proxy $\hat{\mu}_{WHT}$ and $\text{supp}_\varepsilon$.
  \item The remaining $(1-\beta)N$ samples form the validation split used
  to compute the normalized off-attractor mass trajectory and WAC products
  on out-of-sample data.
\end{itemize}

\section{Governed Retry and CAS Commitment}

\subsection{Two-Stage CAS Commitment}

Define a base parameter bundle $\Theta_{\text{Base}}$ including
$\text{retry\_nonce}$. The CAS commitment is computed in two stages:
\begin{enumerate}[noitemsep]
  \item Fix $\Theta_{\text{Base}}$ with $\text{retry\_nonce} = 0$.
  \item Derive $\Theta_{C6}$ and $\text{shuffle\_seed}$ via
  \[
    \text{shuffle\_seed} =
    \text{Poseidon2}(\Gamma_d \Vert \Theta_{\text{Base}}),
  \]
  and then assemble the full $\Theta_{C6}$.
\end{enumerate}
Retries increment $\text{retry\_nonce}$ and therefore yield distinct
CAS commitments, making re-shuffles a visible ACE budget burn.

\section{Implementation Snippets}

\subsection{Python: ThetaC6 and CAS Registry}

\begin{lstlisting}[language=Python,caption={ThetaC6 and CAS registry}]
from dataclasses import dataclass
from typing import Any
from poseidon2 import poseidon_hash

@dataclass(frozen=True)
class ThetaBase:
    epsilon: float
    delta: float
    N0: int
    K: int
    M: int
    beta: float
    tau_min: float
    alpha_M: float
    retry_nonce: int
    domain: str

@dataclass(frozen=True)
class ThetaC6:
    epsilon: float
    delta: float
    N0: int
    K: int
    M: int
    beta: float
    tau_min: float
    alpha_M: float
    shuffle_seed: int
    domain: str
    retry_nonce: int

def derive_cas_registry(gamma_d: bytes, theta_base: ThetaBase) -> ThetaC6:
    # Derive shuffle_seed deterministically from CAS gamma_d and ThetaBase
    seed_bytes = poseidon_hash(gamma_d + serialize(theta_base))
    shuffle_seed = int.from_bytes(seed_bytes[:8], "big")
    return ThetaC6(
        epsilon=theta_base.epsilon,
        delta=theta_base.delta,
        N0=theta_base.N0,
        K=theta_base.K,
        M=theta_base.M,
        beta=theta_base.beta,
        tau_min=theta_base.tau_min,
        alpha_M=theta_base.alpha_M,
        shuffle_seed=shuffle_seed,
        domain=theta_base.domain,
        retry_nonce=theta_base.retry_nonce,
    )
\end{lstlisting}

\subsection{Python: WAC Window Product and L\_s Monitoring}

\begin{lstlisting}[language=Python,caption={WAC window product computation}]
import numpy as np

def compute_wac_window(m_bar: np.ndarray, t: int, m: int) -> float:
    # Windowed average contraction product over [t, t+m]
    num = m_bar[t + m]
    den = max(m_bar[t], 1e-12)
    return num / den
\end{lstlisting}

\subsection{Python: Trait Predicate R1CS Encoding}

\begin{lstlisting}[language=Python,caption={Trait predicate as R1CS constraint}]
def trait_predicate_constraints(h_hat_zeta: int,
                                h_hat_dc: int,
                                theta_trait: int,
                                field_modulus: int):
    # Enforce: h_hat(zeta_trait) >= theta_trait * N, where N = h_hat[0]
    # In R1CS: (theta * N) -> mul gate, then range / comparison.
    thetaN = (theta_trait * h_hat_dc) % field_modulus
    assert h_hat_zeta >= thetaN, "Trait predicate violated"
\end{lstlisting}

\section{7-Day Implementation Plan}

\subsection{Day 1: CAS Registry \& Spec}

\begin{itemize}[noitemsep]
  \item Implement \texttt{ThetaBase}, \texttt{ThetaC6}, and
  \texttt{derive\_cas\_registry} in \texttt{pirtm/cas\_registry.py}.
  \item Update the mathematical specification (this report) to include
  \texttt{retry\_nonce}, $\alpha_M$, and two-stage CAS commitment.
\end{itemize}

\subsection{Day 2: Audit Whitelist \& Constraint Budget}

\begin{itemize}[noitemsep]
  \item Implement \texttt{\_IFMD\_AUDIT\_WHITELIST} and
  \texttt{WitnessFieldRegistry.check} in \texttt{pirtm/audit.py}.
  \item Lock the 5087-constraint canonical budget in \texttt{pirtm/csc.py}.
\end{itemize}

\subsection{Days 3--7: Spectral Governor, Recurrence, and Track A}

\begin{itemize}[noitemsep]
  \item Implement \texttt{govern\_accumulator} to compute $X_n$, SNR, and support size.
  \item Extend \texttt{run()} to use \texttt{shuffle\_seed}, proxy distribution,
  and holdout split for WAC.
  \item Execute the three-split E. coli K-12 experiment and record the final
  JSON artifact for CIP-1 filing.
\end{itemize}

\section{Conclusion}

The progression from the original Codon-Contrast encoding through FV-6 to
the IFMD-grounded FV-8 architecture yields a fully governed, mathematically
sound and ZK-realizable spectral framework. The 6-bit biophysical mapping and
length-64 FWHT are now embedded in an ACE-compliant, PIRTM-governed runtime
with explicit retry accounting and a deterministic contraction certificate
for genomic codon populations.

\nocite{*}
\bibliographystyle{plainnat}
\bibliography{references} 
\end{document}